\section{Signed rough static potentials and the physical wave}
\label{S:signed-static-wave}

Fix the original nonzero covariance and a finite physical mask $m_R$ as in
Section~\ref{S:actual-cameron-mean}. For real $q \in C_c^\infty(\R^d)$ put $V_q
= m_R(\Gamma*q)$ and $L_q(u,v,c_1,c_2) = (0,V_q(u+c_1),0,0)$ on $F =
E\oplus\R^2$. Let $\mathcal{V}_R$ be the operator-norm closure of these
operators in $\mathcal{K}(F)$. It is a closed real linear space. No positivity
is imposed on $q$.

\begin{lemma}[The limiting multiplier]\label{L:signed-limiting-multiplier}
  Every $L \in \mathcal{V}_R$ determines a unique real $V \in H^{-1/2}$
  supported in $\operatorname{supp}m_R$. It has the multiplier action
  $L(u,v,c_1,c_2) = (0,Vu+c_1V,0,0)$, where multiplication extends from tests to
  $H^{1/2} \to H^{-1/2}$. For $k = \|L\|$, its symmetric pairing obeys
  \begin{equation}\label{E:signed-multiplier-form}
    |v_V(u,w)|\le k\|u\|_{H^{1/2}}\|w\|_{H^{1/2}} .
  \end{equation}
\end{lemma}

\begin{proof}
  Choose $L_{q_n} \to L$. The constant-displacement vector has norm one, so
  $\|V_{q_n}-V_{q_j}\|_{H^{-1/2}} \le \|L_{q_n}-L_{q_j}\|$. Thus, $V_{q_n} \to
  V$ in $H^{-1/2}$. Testing outside the fixed compact mask support preserves
  that support. Every zero column of $L_{q_n}$ remains zero. Its displacement
  multipliers converge in $\mathcal{L}(H^{1/2},H^{-1/2})$ to a map of norm at
  most $k$. For compact smooth $u$, smooth multiplication is bounded on
  $H^{-1/2}$ by Lemma~\ref{L:cutoff-factorization}. Hence, $V_{q_n}u \to Vu$
  there, identifying the limiting map on tests and then uniquely on $H^{1/2}$ by
  density. The constant column recovers $V$ and proves uniqueness. On compact
  smooth $u,w$, the pairing is $\langle V,uw\rangle$, denoted $v_V(u,w)$, so it
  is symmetric. Its bound and symmetry extend by density. Complexification
  preserves the operator norm by splitting real and imaginary parts in the
  Hilbert norms. The pairing extends sesquilinearly to a Hermitian form; its
  diagonal is the sum of the real and imaginary diagonal values. This proves
  Lemma~\ref{L:signed-limiting-multiplier}.
\end{proof}

The bound~\eqref{E:signed-multiplier-form} controls the pairing $v_V$ on
$H^{1/2}$. We now use it to define a self-adjoint operator $H_V$ from the
quadratic form $q_V$ on $H^1$. We also compare the quarter power of its shift
$A$ with the $H^{1/2}$ norm.

\LeanMarker{P:signed-closed-form}
\begin{proposition}[Closed form and quarter-power phase]
  \label{P:signed-closed-form}
  Let $K \ge \max(1,\|L\|)$ and $c = 2+K^2$. On complexified $L^2$ the form
  \[
    q_V(u)=\|\nabla u\|_2^2-v_V(u,u),\qquad Q(q_V)=H^1,
  \]
  represents a unique self-adjoint $H_V$. With $A = H_V+c$ and $L_K = 3(1+K^2)$,
  \begin{align}
    A&\ge1,\qquad D(A^{1/2})=H^1,\label{E:signed-form-domain}\\
    \tfrac12\|u\|_{H^1}^2
    &\le q_V(u)+c\|u\|_2^2\le L_K\|u\|_{H^1}^2,\label{E:signed-form-bounds}\\
    D(A^{1/4})&=H^{1/2},\qquad
    2^{-1/4}\|u\|_{H^{1/2}}
    \le\|A^{1/4}u\|_2\le L_K^{1/4}\|u\|_{H^{1/2}}.
    \label{E:signed-quarter-phase}
  \end{align}
  The negative dual completion is $H^{-1/2}$. No $H^2$ operator domain is
  asserted.
\end{proposition}

\begin{proof}
  Put $b = \|u\|_2$, $z = \|u\|_{H^1}$. Fourier Cauchy--Schwarz
  and~\eqref{E:signed-multiplier-form} give
  \[
    \|u\|_{H^{1/2}}^2\le bz,\qquad
    |v_V(u,u)|\le Kbz\le\tfrac12z^2+\tfrac12K^2b^2 .
  \]
  The last scalar inequality is the weighted arithmetic--geometric mean
  inequality at weights $1/2,1/2$ and nonnegative bases $z^2,(Kb)^2$, including
  zero norms. Since $q_V = z^2-b^2-v_V(u,u)$, it follows that
  \[
    q_V+cb^2\ge\tfrac12z^2+(1+K^2/2)b^2,\qquad
    q_V+cb^2\le\tfrac32z^2+(1+3K^2/2)b^2\le L_Kz^2.
  \]
  The lower bound dominates both $z^2/2$ and $b^2$. Thus, this densely defined
  Hermitian form is semibounded and its shifted norm is equivalent to the
  complete $H^1$ norm. It is closed in the definition
  of~\cite{teschl:14:mathematical}*{p.~80}.
  From~\cite{teschl:14:mathematical}*{Theorem~2.14, p.~81, (2.63)}, it
  represents a unique self-adjoint $H_V$. This representation is independent of
  the admissible shift. From that book's Theorem~3.6 and~(3.51)--(3.53), p.~109,
  applied after shifting, we obtain~\eqref{E:signed-form-domain}. Here $A \ge 1$
  is initially its form lower bound. A bounded spectral slice below one would
  produce a unit $D(A)$ test contradicting that form bound; hence it is also the
  spectral lower bound. Lemma~\ref{L:signed-spectral-wave} records the full
  bottom argument.

  Set $B_0 = 1-\Delta$. On the common form domain, $\tfrac12B_0 \le A \le
  L_KB_0$. We prove the needed square-root comparison rather than assuming an
  order rule for arbitrary operator powers. For $T_0,S_0 \ge 1$ with common form
  domain and $q_{T_0} \le q_{S_0}$, completing the square in the shifted form
  norm gives
  \[
    \langle f,(T_0+t)^{-1}f\rangle
    =\sup_u\{2\operatorname{Re}\langle f,u\rangle
    -q_{T_0}(u)-t\|u\|^2\}
    \ge\langle f,(S_0+t)^{-1}f\rangle,\quad t>0 .
  \]
  The scalar identity $\sqrt x = \pi^{-1}\int_0^\infty t^{-1/2}x/(x+t)\ud t$
  follows by $t = xv^2$. Positive spectral integration and Tonelli now imply
  $T_0^{1/2} \le S_0^{1/2}$ as forms, with their domain inclusions:
  $I-t(T_0+t)^{-1} \le I-t(S_0+t)^{-1}$. Apply this to $(B_0,2A)$ and
  $(A,L_KB_0)$. Both operators in each pair are at least one. The resulting
  inequalities give~\eqref{E:signed-quarter-phase}, including equality of its
  domains. For $v \in L^2$ the $L^2$-duality norm against $\|A^{1/4}u\|_2$ is
  $\|A^{-1/4}v\|_2$: Cauchy--Schwarz gives the upper bound, and $u = A^{-1/2}v$
  gives equality after normalization. Duality and completion therefore identify
  the negative phase. Uniqueness of the form, resolvent and spectral calculus
  makes these constructions commute with real conjugation. This proves
  Proposition~\ref{P:signed-closed-form}.
\end{proof}

Proposition~\ref{P:signed-closed-form} gives the operator $H_V$ and compares
the norms built from $A^{1/4}$ with those of $E$. We now locate the bottom of
the spectrum of $H_V$ and bound the wave that its spectral calculus defines on
$E$. The bound allows exponential growth at the rate $\lambda(V)$.

\begin{lemma}[The bottom and spectral wave]\label{L:signed-spectral-wave}
  Write $P_V^E$ for the evolution in~\eqref{E:signed-spectral-matrix}, and $S_E$
  for the free restriction from Lemma~\ref{L:constant-phase-free}. Put
  \[
    \beta(V)=\inf_{\substack{u\in H^1\\\|u\|_2=1}}q_V(u),
    \qquad \lambda(V)=\sqrt{-\beta(V)} .
  \]
  Then $\beta(V) = \inf\operatorname{spec}H_V$ and $-K^2 \le \beta(V) \le 0$.
  The spectral wave on $E$ has the bound
  \begin{equation}\label{E:signed-E-wave}
    \|P_V^E(t)\|\le8(1+K)^2(1+t)e^{\lambda(V)t},\qquad t\ge0.
  \end{equation}
\end{lemma}

\begin{proof}
  For unit $u$, $z = \|u\|_{H^1} \ge 1$ gives $q_V(u) \ge z^2/2-1-K^2/2 \ge
  -(1+K^2)/2 \ge -K^2$. Conversely dilate a unit compact smooth bump and
  translate its whole support outside $\operatorname{supp}m_R$. Its potential
  pairing is zero and its kinetic energy tends to zero. Thus, $\beta \le 0$.
  Real and complex form infima agree because a complex diagonal splits into the
  two real diagonals. The spectral form formula
  in~\cite{teschl:14:mathematical}*{(3.52)--(3.53)} identifies this form infimum
  with the spectral bottom. For clarity, a bounded spectral slice below the form
  lower bound would give a unit form test violating that bound. Conversely the
  form integral is bounded below by the spectral infimum. The spectrum is
  nonempty: a unit vector's scalar spectral measure has mass one and a support
  point, and normalized projections onto arbitrarily small bounded neighborhoods
  give approximate eigenvectors, excluding a bounded inverse there. Conversely,
  off the closed PVM support $\Sigma$ the Borel function $(h-z)^{-1}$, set to
  zero off $\Sigma$, is bounded for real $z\notin\Sigma$, as is $h/(h-z)$ on
  $\Sigma$. Its calculus maps into $D(H_V)$ and is a two-sided resolvent inverse
  on the appropriate domains. Thus, the spectrum is precisely $\Sigma$. Hence,
  $\lambda \le K$ and every spectral value $h$ has $h+c \ge 2$.

  Use the phase norm $\|(u,v)\|_{E_A}^2 = \|A^{1/4}u\|_2^2+\|A^{-1/4}v\|_2^2$.
  With $M_K = L_K^{1/4}$, the preceding proposition and its dual give
  $M_K^{-1}\|z\|_E \le \|z\|_{E_A} \le M_K\|z\|_E$. The condition-number cost is
  $M_K^2 = \sqrt{3(1+K^2)} \le 2(1+K)$. Define $C_t(h),S_t(h)$ by $\cos(t\sqrt
  h),\sin(t\sqrt h)/\sqrt h$ for $h > 0$, by $1,t$ at zero, and by
  $\cosh(t\sqrt{-h}),\sinh(t\sqrt{-h})/\sqrt{-h}$ for $h < 0$. In the $E_A$
  isometry the wave matrix is
  \begin{equation}\label{E:signed-spectral-matrix}
    \begin{pmatrix}
      C_t(h)&\sqrt{h+c}\,S_t(h)\\
      -hS_t(h)/\sqrt{h+c}&C_t(h)
    \end{pmatrix}.
  \end{equation}
  For $h \ge 0$, the entries are bounded using $|C_t| \le 1$, $|S_t| \le
  \min(t,h^{-1/2})$, $\sqrt{h+c}|S_t| \le 1+\sqrt c\,t$ and $|hS_t|/\sqrt{h+c}
  \le 1$; at zero the off-diagonals are $\sqrt c\,t,0$. For $\beta \le h < 0$,
  they are bounded using $|C_t| \le e^{\lambda t}$, $|S_t| \le te^{\lambda t}$,
  $\sqrt{h+c}|S_t| \le \sqrt c\,te^{\lambda t}$ and $|hS_t|/\sqrt{h+c} \le
  Ke^{\lambda t}$. There is no division by $\lambda$. Since $\sqrt c \le 2K$,
  each entry is at most $2(1+K)(1+t)e^{\lambda t}$. Bounded Borel calculus
  from~\cite{teschl:14:mathematical}*{Theorem~3.1, pp.~102--103,
  (3.18)--(3.21)}, the two-by-two matrix norm bound and the condition-number
  cost give~\eqref{E:signed-E-wave}. Functions can be set to zero off the
  spectral support. Spectral dominated convergence gives strong continuity, and
  scalar addition identities give the group property. The same matrix and bounds
  apply on the weighted graph levels $D(A^{1/2})\times L^2$ and $D(A)\times
  D(A^{1/2})$. Spectral cutoffs justify their extensions and continuity there.
  This proves Lemma~\ref{L:signed-spectral-wave}.
\end{proof}

Lemma~\ref{L:signed-spectral-wave} bounds a wave defined by spectral
calculus, and only on $E$. We still have to identify it with the physical wave
and to include the two constant columns of $F$.

\begin{proposition}[Identification with the physical wave]
  \label{P:signed-physical-wave}
  The spectral wave in Lemma~\ref{L:signed-spectral-wave} is the
  bounded-perturbation, equivalently Young, physical wave on $E$. Its extension
  on $F$ satisfies
  \begin{equation}\label{E:signed-F-wave}
    \|P_V(t)\|_{\mathcal{L}(F)}
    \le16(1+K)^4(1+t)^3e^{\lambda(V)t},\qquad t\ge0.
  \end{equation}
  Both constant response columns are included.
\end{proposition}

\begin{proof}
  The imported form representation \cite{teschl:14:mathematical}*{Theorem~2.14,
  p.~81, (2.63)}, tested against compact smooth functions, gives $H_Vu = -\Delta
  u-Vu$ in distributions for $u \in D(H_V)$. Such $u$ lies in $H^1$ and $Vu \in
  H^{-1/2}$. Consequently, $(1-\Delta)u = H_Vu+Vu+u \in H^{-1/2}$, so
  $D(H_V)\hookrightarrow H^{3/2}$. The embedding is graph-continuous:
  by~\eqref{E:signed-form-bounds}, $\|u\|_{H^1}^2 \le 2\langle Au,u\rangle \le
  \|Au\|_2^2+\|u\|_2^2$, and the distributional equation bounds the $H^{3/2}$
  norm by this graph norm.

  The generator of the free $E$ group is $A_0(u,v) = (v,\Delta u)$ on
  $H^{3/2}\times H^{1/2}$. Indeed, differentiate its Fourier matrix after a
  frequency cutoff, dominate by that graph norm and remove the cutoff.
  Conversely any strong generator derivative has those distributional
  components; they require $v \in H^{1/2}$ and $\Delta u \in H^{-1/2}$, hence $u
  \in H^{3/2}$. No operator-norm differentiability of the group is used.

  Start spectral data in $D(A)\times D(A^{1/2}) = D(H_V)\times H^1$. The
  preceding graph-level estimates give continuous graph trajectories and $C^1$
  trajectories in $H^1\times L^2$, with $u_t = v$, $v_t = -H_Vu$. This follows
  first on bounded spectral cutoffs and then uniformly on finite horizons in the
  graph norms by spectral dominated convergence
  and~\cite{teschl:14:mathematical}*{Theorem~3.2, p.~104, (3.29)--(3.33)}. The
  projections $1_{[1,n]}(A)$ approximate every graph datum and the same weighted
  matrix bounds control all cutoff tails on those horizons. Graph embedding
  gives continuity in $D(A_0)$ and $z_t = A_0z+L_Ez$ in $E$, where $L_E = L|_E$.
  The strong-group chain rule on this domain yields
  \[
    S_E(-t)z(t)=z(0)+\int_0^t S_E(-s)L_Ez(s)\ud s .
  \]
  Graph data are dense in $E$: spectral cutoffs make $D(A)$ dense in $D(A^{1/2})
  = H^1$, and $H^1$ is dense in both $H^{1/2}$ and $H^{-1/2}$. Boundedness on
  finite horizons extends this identity to all $E$. Volterra uniqueness by
  Picard iteration identifies the physical wave. No $D(A^{3/4}) = H^{3/2}$
  assertion is needed.

  On $F$, bounded $L$ and strongly continuous $S_F$ similarly define a unique
  Duhamel flow. On $[0,T]$, the $n$th Volterra term has bound
  $C_T^{n+1}\|L\|^nT^n/n!$. Uniqueness makes horizons consistent and gives
  composition. Its interaction primitive is $\int_0^t S_F(-s)LS_F(s)\ud s$.
  Compactness of $L$, and strong continuity of the group and its adjoint, give
  norm-continuity of the conjugated compact operator, as in
  Proposition~\ref{P:cameron-actual-embedding}. Thus, this primitive is
  Lipschitz and its ordinary equation is exactly the Young equation.

  For initial $(e,c_1,c_2)$ the constant coordinates are $(c_1+tc_2,c_2)$ and
  the $E$ component is
  \[
    P_V^E(t)e+\int_0^t P_V^E(t-s)(0,V(c_1+sc_2))\ud s .
  \]
  The bounded $E$-valued Fubini argument proving \eqref{E:compact-F-block}
  applies: substitute the $E$ Duhamel formula under this integral and
  interchange the absolutely bounded finite time triangle. It solves the full
  system, and uniqueness identifies it. No ground positivity or gap is used. The
  forcing norm is at most $K(1+s)|(c_1,c_2)|$. Since $\lambda \ge 0$ and
  $\int_0^t(1+t-s)(1+s)\ud s \le (1+t)^3$, the full norm is at most
  \[
    (1+t)^3e^{\lambda t}
    \{8(1+K)^2\|e\|+[8K(1+K)^2+1]|(c_1,c_2)|\}.
  \]
  Scalar Cauchy--Schwarz and $\sqrt{8^2+9^2} < 13$ bound this by
  $13(1+K)^3(1+t)^3e^{\lambda t}\|(e,c_1,c_2)\|_F$. This implies the stated
  sufficient bound. This proves Proposition~\ref{P:signed-physical-wave}.
\end{proof}
